\pagestyle{empty}

\noindent \textbf{\Large A CD-melléklet tartalma}

\vskip 1cm

\noindent A dolgozathoz mellékelt lemezen a \texttt{diplomamunka} nevű könyvtárban az alábbi fájlok találhatók:

\begin{itemize}
\item a dolgozat \LaTeX\ forráskódja,
\item a dolgozat PDF formátumban (\texttt{dolgozat.pdf}).
\item A dolgozat kiírása Word formátumban a \texttt{docs} jegyzékben \\ (\texttt{Drig\_Dávid\_diplomamunka\_kiiras.docx}).
\item A magyar és angol nyelvű összefoglaló \LaTeX\ és PDF formátumban \\ (\texttt{8\_osszegzes.tex}, \texttt{8\_osszegzes.pdf}, \texttt{9\_summary.tex}, \texttt{9\_summary.pdf}).
\end{itemize}

A \texttt{doc} nevű könyvtár tartalmazza a dolgozathoz kapcsolódó dokumentációkat, valamint a Sphinx alapú \texttt{github.io} oldalt.

A \texttt{notebooks} könyvtárban a dolgozat elkészítése során létrehozott, illetve a dolgozatban ismertetett Jupyter-munkafüzetek, a hozzájuk tartozó Python-szkriptek, továbbá a munkafüzetek HTML formátumú megjelenítéséhez szükséges szkript és maguk a HTML formátumú munkafüzetek találhatók.
